% Cập nhật ngày 15/10/2026. Nguồn mới: results/stage1b-report.txt, mục 2, 5.
\chapter*{Tóm tắt}
\addcontentsline{toc}{chapter}{Tóm tắt}

Trong dịch vụ đi chung xe, giờ trả đã hứa có thể bị dời khi bộ điều phối ghép thêm khách hoặc khi thời gian đi thay
đổi. Khóa luận đặc tả và cài đặt luật C trong RideBound: tổng độ lớn phần dời giờ trả tính cho điều phối của mỗi
khách trong cả chuyến không vượt một trần cố định. Mỗi lần dời được đo so với giờ trả tính lại khi giữ nguyên kế
hoạch cũ, mốc đã có ở Hörl và Zwick. Luật C đi cùng lời hứa theo dự báo tắc đường và sổ theo từng khách; sổ phân
biệt phần tính cho điều phối, phần được dự báo giải thích và phần còn lại.
% EV: A01; A06; X3-12; J01; J03; J35. Giới hạn có điều kiện ở muc:giunguyen.

Phép đo dùng FleetPy trên 34 kịch bản Manhattan, với 3.672 yêu cầu mỗi cấu hình. Khóa luận so hai cách chọn phương
án: ưu tiên ổn định và ưu tiên chi phí kế hoạch. Khi cùng giữ trần 30 giây, cách thứ hai phục vụ 2.364 thay vì 2.338
lượt và có 432 thay vì 313 khách hoàn thành chuyến có đi chung xe. Đổi lại, số khách bị sửa giờ trả tăng từ 33 lên
106. Không khách nào vượt trần ở cả hai cách chọn. Với cách ưu tiên chi phí, bỏ trần tăng số lượt phục vụ từ 2.364
lên 2.502, nhưng số lượt trễ hơn lời hứa đầu chưa tính dự báo quá một phút tăng từ 72 lên 379. Các nhóm khách được
phục vụ khác nhau, nên các chênh lệch mô tả đánh đổi giữa các cấu hình.
% EV: S15-01 (report mục 2, 5d--5f); S15-02 (kiểm lại số từ trích xuất).

Khi tính lại lời hứa có dự báo trên hành trình đã ghi của cấu hình ưu tiên chi phí có trần, số lượt trễ quá một
phút là 11 trong 2.364 lượt phục vụ. Đây là điều chỉnh giờ hứa, không làm xe chạy nhanh hơn. Đối chứng hạn chót
và đệm lời hứa cho thấy luật C không vượt trội ở mọi thước đo. Các kết quả có tính thăm dò, chỉ áp cho mô phỏng tải
thấp với hệ số tắc đường chung toàn thành phố. Khóa luận chưa chọn một trần tối ưu hoặc xác nhận mức hài lòng của
khách; phần dự báo dùng hồ sơ lịch sử đơn giản, chưa phải một mô hình dự báo giao thông mới.
% EV: S15-01 (5b); C01; C10; K05; K36; I01; I05; F01--F04.

\vspace{6pt}
\noindent\textbf{Từ khóa:} điều phối xe đi chung, lời hứa giờ trả, trần theo từng khách, lời hứa theo dự báo,
sổ theo từng khách, mô phỏng FleetPy.
\clearpage

\chapter*{Abstract}
\addcontentsline{toc}{chapter}{Abstract}

In ride-pooling, promised drop-off times may change when a dispatcher inserts riders or travel times change.
This thesis specifies and implements rule C in RideBound: each rider has a fixed cap on the sum of the magnitudes
of dispatch-charged drop-off changes over the trip. Each change is measured against the updated drop-off time
under the existing plan, a reference already used by Hörl and Zwick. Rule C is combined with forecast-based
promises and a per-rider ledger separating dispatch-charged changes, the part explained by the forecast, and the
remainder.

Measurements use FleetPy on 34 Manhattan scenarios, with 3,672 requests per configuration. Two selection orders
are compared: stability-first and planned-cost-first. Under the same 30-second cap, planned-cost-first serves
2,364 rather than 2,338 riders, and 432 rather than 313 served riders share a vehicle. In exchange, the number
whose drop-off promise is revised rises from 33 to 106. Neither capped configuration exceeds its dispatch-change
budget. Under planned-cost-first, removing the cap raises service from 2,364 to 2,502, but arrivals more than one
minute after the initial unadjusted promise rise from 72 to 379. Served populations differ, so these are
descriptive configuration trade-offs.

Replaying forecast-based promises on the capped planned-cost-first trajectories leaves 11 of 2,364 served
riders more than one minute late. This changes the quoted time, not vehicle movement. Deadline and padding
comparisons show that rule C does not dominate every measure. The findings are exploratory and limited to the
low-load simulation with city-wide traffic factors. No optimal budget or passenger-satisfaction claim is made;
the forecast uses a simple historical profile rather than a new traffic-prediction model.

\vspace{6pt}
\noindent\textbf{Keywords:} ride-pooling dispatch, promised drop-off time, per-rider cap, forecast-based promise,
per-rider ledger, FleetPy simulation.
\clearpage
